\documentclass[11pt,a4paper]{article}
\usepackage[margin=2.2cm]{geometry}
\usepackage{lmodern}
\usepackage[T1]{fontenc}
\usepackage{listings}
\usepackage{xcolor}
\usepackage{hyperref}
\hypersetup{colorlinks=true,linkcolor=blue,urlcolor=blue}

\lstset{
  basicstyle=\ttfamily\small,
  frame=single,
  framerule=0.3pt,
  rulecolor=\color{black!30},
  backgroundcolor=\color{black!3},
  showstringspaces=false,
  breaklines=true,
  columns=fullflexible,
  literate={~}{{\textasciitilde}}1,
  xleftmargin=0.5em,
  xrightmargin=0.5em,
}

\title{\textbf{dyablo FoF: user guide}}
\author{}
\date{}

\begin{document}
\maketitle
\vspace{-2em}

\noindent
A parallel Friends-of-Friends halo finder for dyablo HDF5 particle snapshots.
Adapted from the RAMSES-era pFoF (Roy/Audit) tool. This guide covers building,
running, and post-processing the output to produce a halo mass function plot
against theory.

\section{What you get}

The shipped directory contains:
\begin{itemize}\itemsep0pt
  \item \texttt{*.f90, makefile} --- the Fortran sources and build system.
  \item \texttt{fof.in} --- example input file (one line per parameter).
  \item \texttt{plot\_hmf.py} --- Python script that reads the FoF catalog and
        overplots Sheth--Tormen / Tinker08 at the snapshot redshift.
\end{itemize}

\section{Build}

Requirements: an MPI-aware Fortran compiler (\texttt{mpif90}) and HDF5 with
Fortran bindings.

The \texttt{makefile} is set up for macOS with Homebrew (\texttt{brew install
hdf5 open-mpi}). The only two lines you typically have to touch are at the
top:

\begin{lstlisting}
HDF5_DIR ?= /opt/homebrew                 # <-- path to HDF5 install prefix
                                          #     (must contain include/ and lib/)

OPTS ?= -O3 -cpp -DLONGINT \              # -DLONGINT activates 8-byte particle
        -I$(HDF5_DIR)/include \           # IDs; required for N_p > 2^31.
        -ffree-line-length-none           # GFortran-specific; remove for ifort.

LIBS ?= -L$(HDF5_DIR)/lib -lhdf5_fortran -lhdf5
\end{lstlisting}

\noindent
Override on the command line if needed:
\begin{lstlisting}
make HDF5_DIR=/usr/local         # custom HDF5 prefix
make MPIFC=mpiifort FC=mpiifort  # Intel MPI / ifort
\end{lstlisting}

Then simply:
\begin{lstlisting}
make
\end{lstlisting}
produces the \texttt{fof} executable. \texttt{make clean} removes objects and the binary.

\section{The fof.in input file}

\texttt{fof.in} is read line by line --- order matters; inline comments are
allowed only on numerical/logical lines (lines read with list-directed input),
not on string lines.

\begin{lstlisting}
my_halos                  ! root: output filename prefix
DY                        ! code_index (must be 'DY' for dyablo)
/path/to/snapshot/dir     ! pathinput: dir containing the .h5 and .ini
star_particles_..._iter0070000.h5   ! particle HDF5 file
restart_star_NNNNNN.ini             ! the dyablo restart .ini
0       !grpsize (RAMSES legacy, leave 0)
0.2     !percolation parameter b (in mean-inter-particle units)
100     !min mass (in particles): halos with fewer particles are not written
100000000 !max mass (in particles)
.true.  !star: snapshot contains star particles to filter out
.false. !metal: never .true. for dyablo
.false. !outcube: write per-rank particle cubes after read (debug)
.true.  !dofof: actually run halo detection (set .false. to just read+exit)
.false. !readfromcube: read particles from precomputed cubes (debug)
.true.  !dotimings: extra MPI syncs to print accurate timings
.false. !invent_ids: force-invent IDs even when snapshot has /id (legacy)
outputs                   ! outdir: where to write outputs (created if missing)
\end{lstlisting}

\noindent
Notes:
\begin{itemize}\itemsep0pt
  \item \texttt{root}, \texttt{pathinput}, \texttt{namepart}, \texttt{nameinfo},
        and \texttt{outdir} are read as 80-char strings; do not append comments
        on those lines.
  \item DM particle mass is \emph{auto-detected} as the median of the
        \texttt{mass} dataset, so you don't have to set it.
  \item By default the code uses the original \texttt{/id} dataset if present.
        Set \texttt{invent\_ids = .true.} only if your snapshot has duplicate
        or corrupt IDs.
\end{itemize}

\section{Run}

\begin{lstlisting}
mpirun -np 1 ./fof          # serial reference
mpirun -np 8 ./fof          # 2 x 2 x 2 spatial decomposition
mpirun -np 64 ./fof         # 4 x 4 x 4
\end{lstlisting}

\noindent\textbf{Rank-count constraint.} The MPI decomposition is a cubic
$n_{\mathrm{sd}}\times n_{\mathrm{sd}}\times n_{\mathrm{sd}}$ grid with
$n_{\mathrm{sd}}=\lfloor N_{\mathrm{proc}}^{1/3}\rfloor$, and the particle grid
$n_{\mathrm{res}}$ (from \texttt{particle\_grid.nx} in the \texttt{.ini}) must
be divisible by $n_{\mathrm{sd}}$. For $n_{\mathrm{res}}=256$ the valid rank
counts are therefore $1, 8, 64, 512, \ldots$ (i.e.\ cubes of powers of two).

\section{Output files}

In \texttt{outdir/}:
\begin{itemize}\itemsep0pt
  \item \texttt{root.log} --- run log (parameters, timings, halo count).
  \item \texttt{root.inp} --- echo of the input parameters.
  \item \texttt{root\_masst\_NNNNN} --- one file per MPI rank: halo catalog
        (Fortran sequential unformatted). Each halo record is
        \texttt{int64 id, int32 mass, 3*float32 (x,y,z)} with positions
        normalized to $[0,1]^3$.
  \item \texttt{root\_strct\_NNNNN} --- per-halo particle membership (positions,
        velocities, and original snapshot IDs).
\end{itemize}

\section{Plotting the halo mass function}

\texttt{plot\_hmf.py} reads the \texttt{*\_masst\_*} files, pulls the cosmology
from the dyablo \texttt{.ini} and the current scale factor from the snapshot
\texttt{.h5} (\texttt{scalar\_data/aexp}), and overplots two analytic HMFs
(Sheth--Tormen and Tinker08) at the same redshift.

\subsection*{Quick start: reuse your \texttt{fof.in}}

The simplest invocation just points the script at the same \texttt{fof.in}
you ran the finder with:

\begin{lstlisting}
python3 plot_hmf.py --fofin fof.in
\end{lstlisting}

\noindent From \texttt{fof.in} the script derives every input it needs, so you
do not have to retype the paths:
\begin{itemize}
  \item \texttt{--catalog} from the \texttt{outdir} field (resolved relative to
        the \texttt{fof.in} directory, as the finder itself does);
  \item \texttt{--ini} from the \texttt{nameinfo} field, under \texttt{pathinput};
  \item \texttt{--particle-h5} from the \texttt{namepart} field;
  \item \texttt{--h5} (the state file carrying \texttt{scalar\_data/aexp}) from
        the same particle name with dyablo's \texttt{\_particles\_particles}
        infix stripped (e.g.\ \texttt{star\_particles\_particles\_iterN.h5}
        $\rightarrow$ \texttt{star\_iterN.h5});
  \item \texttt{--b} (linking length) from the \texttt{perco} field.
\end{itemize}

\noindent Any of these given explicitly on the command line override the value
taken from \texttt{fof.in}, so you can mix the two --- e.g.\
\texttt{plot\_hmf.py --fofin fof.in --out z5.png}.

\subsection*{Or pass the paths explicitly}

\begin{lstlisting}
python3 plot_hmf.py \
    --catalog outputs/ \
    --ini     /path/to/restart_star_NNNNNN.ini \
    --h5      /path/to/star_iterNNNNNNN.h5
\end{lstlisting}

\noindent Optional flags: \texttt{--sigma8} (default 0.81), \texttt{--ns}
(default 0.96), \texttt{--hmf-model} (default \texttt{SMT}; any model name
accepted by the \texttt{hmf} package), \texttt{--out} (output PNG path,
default \texttt{hmf\_compare.png}).

\noindent Python dependencies: \texttt{numpy scipy h5py matplotlib hmf astropy}.

\noindent \emph{Note:} $\sigma_8$ and $n_s$ are not stored in the dyablo
\texttt{.ini}; the script uses Planck-ish defaults. Override with
\texttt{--sigma8/--ns} if your ICs use different values.

\section{Troubleshooting}

\begin{itemize}\itemsep0pt
  \item \emph{Build fails on \texttt{Cannot open .mod}}: HDF5 Fortran module
        not found. Check \texttt{\$(HDF5\_DIR)/include}.
  \item \emph{Run dies in \texttt{dyablo\_lecture}}: usually a path issue ---
        check \texttt{pathinput}, \texttt{namepart}, \texttt{nameinfo}.
  \item \emph{Crash in \texttt{MPI\_Finalize} on a laptop} (macOS, OFI
        \texttt{utun*} errors): harmless --- the FoF work is already done and
        output files are written. Set \texttt{FI\_PROVIDER=shm} or disconnect
        the VPN to silence it. Does not occur on HPC clusters.
\end{itemize}

\section{History and acknowledgements}

This tool descends from the parallel Friends-of-Friends code originally written
for RAMSES snapshots: Edouard Audit's sequential FoF (CEA), parallelised by
Fabrice Roy (CNRS / LUTH, Observatoire de Paris), with further contributions
from Yann Razera.

The adaptation to dyablo HDF5 snapshots was developed across two interactive
coding sessions with Claude Opus 4.7 (Anthropic):

\begin{itemize}\itemsep0pt
  \item \emph{Port.} Replaced the RAMSES particle reader with a dyablo HDF5
        reader (\texttt{dyablo\_lecture}); added a minimal parser for the
        dyablo restart \texttt{.ini} to recover box geometry and the particle
        grid resolution; auto-detection of the DM particle mass as the median
        of the \texttt{mass} dataset (so the mass no longer has to be
        hard-coded); optional filtering of star particles via the snapshot's
        \texttt{birth\_time} dataset; preservation of the snapshot's
        \texttt{/id} dataset by default (the legacy ``invent IDs'' behaviour
        is now opt-in via \texttt{invent\_ids = .true.}); removal of RAMSES-only
        code paths.
  \item \emph{Validation and packaging.} Configurable output directory
        (\texttt{outdir}) so runs no longer clutter the source tree;
        cross-rank reproducibility checks (catalog bit-exact at 1 and 8 ranks;
        the few-particle drift seen at 64 ranks is consistent with
        floating-point ambiguity at the larger number of rank boundaries);
        halo-mass-function comparison against Sheth--Tormen and Tinker08, with
        cosmology pulled automatically from the dyablo \texttt{.ini} and the
        current scale factor from \texttt{scalar\_data/aexp}; per-halo
        particle visualisation script for sanity-checking the cosmic web;
        cross-platform build verified on macOS / Homebrew and on Linux /
        Debian (\texttt{h5pfc}); this user guide.
\end{itemize}

\noindent The Fortran sources retain the original authorship attributions at
the top of each module. The dyablo branch was developed and validated with
the assistance of Claude Opus 4.7.

\end{document}
